% Generated from chapter-06.md - do not edit.

\chapter{Better Thoughts Through Better Beliefs}%
\index{better beliefs}\nobreak%

Our thoughts matter not only because of how they make us feel in the moment, but because thoughts repeated over time can grow into beliefs. Those beliefs can become like well-worn ruts in the mind---familiar paths our thoughts return to again and again. But those paths can change. Here is one experience from my own life that showed me how changing my thinking could gradually change what I believed---and, eventually, how I am living.

\section{\emph{How I Learned Food Mastery}}%
\index{food mastery}\nobreak%

\emph{One of the biggest ``better thought'' shifts in my life happened around something very ordinary:}

\textbf{food.}

\emph{For years, I felt out of control with eating.}

\emph{Not wildly so---just enough to complain about it, feel uncomfortable in my clothes, and wish I could do better.}

\emph{I would eat more sugar than I wanted, eat more than I needed, and then feel disappointed with myself.}

\emph{(A very reliable cycle.)}

\emph{I told myself lots of familiar stories:}

\begin{itemize}
\item \emph{It's comforting.}
\item \emph{It tastes so good.}
\item \emph{Holidays and birthdays deserve special extra eating.}
\item \emph{I'll start later.}
\item \emph{The temptation is too strong.}
\item \emph{Everyone else is doing it.}
\end{itemize}

\emph{And underneath all that, there was a belief I hardly noticed at the time:}

\textbf{``I don't really have control here.''}

\section{The Turning Point: A New Belief}

Everything changed when I took Jackie Kelm's class on rapid change.\index{Jackie Kelm}

As part of the class, Jackie had a private conversation with each of us.

She asked me three simple---but powerful---questions.

\section{First:}

\textbf{What is this behavior getting you?}

I told her the truth.

It was comforting.\\
It was fun.\\
It tasted good.\\
It felt like a treat.\\
It felt like I deserved it.

\emph{(A very convincing list.)}

\section{Second:}

\textbf{What is it doing to you?}

That answer was just as honest.

I complained too much.\\
My clothes didn't fit well.\\
I felt out of control.\\
I wasn't happy with myself.

\section{Third:}

\textbf{What do you really want?}

And I heard myself say something new:

\textbf{``I want food mastery.''}\index{food mastery}

\emph{Not dieting.}\\
\emph{Not willpower.}\\
\emph{Not punishment.}

\textbf{Food mastery.}

\emph{I wanted to be able to look at food and calmly decide.}

\emph{To weigh the good feelings of eating it against the downsides.}

\emph{To eat enough so I wasn't truly hungry.}

\emph{To stop feeling driven by temptation.}

\emph{To feel at ease with food.}

\emph{(Imagine that---peace with a cookie.)}

\section{Creating Food Mastery}%
\index{food mastery}\nobreak%

I chose a simple structure: about 1200 calories a day---enough for me to feel satisfied and still move toward my goal weight of 140 pounds.

I started writing my food down.

Not in a harsh way---in a caring, attentive way.

Something surprising happened.

When I was eating enough, I didn't feel desperate.

When I was paying attention, I didn't feel deprived.

And when I saw my progress, I stopped making excuses.

\emph{(My mind had fewer places to hide.)}

\emph{A new belief formed:}

\textbf{I have food mastery.}

\emph{Not I will someday.}\\
Not \emph{I hope to.}\\
Not \emph{I'm trying.}

\textbf{I have it.}

\section{Pause for a Moment}

What is one belief you have about yourself . . . \\
that might not be helping you?

Something like:

\begin{itemize}
\item ``I can't stick with things''
\item ``I don't have discipline''
\item ``This is just the way I am''
\end{itemize}

Now gently wonder:

\textbf{What might be a kinder, truer belief?}

You don't have to force it.

Just consider the possibility.

\section{The Daily Better Thought}

Now, almost every day, I quietly remind myself:

\begin{itemize}
\item I have food mastery.
\item I choose what I eat.
\item I decide what fits in my life.
\item I stay at the weight that feels good to me.
\item I trust myself.
\end{itemize}

When temptation appears, I don't argue with it.

I don't feel bad.

I simply decide.

Sometimes I include the treat.\\
Sometimes I don't.

Either way, I'm in charge.

That is what mastery feels like.

\emph{(Calm . . .  not dramatic.)}

\section{Why This Matters for Better Thoughts}

This whole shift came from changing a belief.

The moment my belief changed from:

\textbf{``Food controls me''}

to:

\textbf{``I have food mastery,''}\index{food mastery}

my thoughts became kinder, calmer, and more capable.

Better beliefs create better thoughts.\\
Better thoughts create better choices.\\
Better choices create a gentler life.

And I no longer plan to ``not care'' about my weight.

I care---kindly, consciously, without drama.

\emph{(Drama is not required for progress.)}

\emph{I care because I love my body.}\\
\emph{I care because I feel better this way.}\\
\emph{I care because I have learned something important about myself:}

\textbf{I can choose.}

\section{Learning Together}

I also know I don't do this alone.

Sometimes my ideas sound a little unusual when I say them out loud.

\emph{(That's often a good sign.)}

\emph{But when I talk with others, something shifts.}

\emph{My thinking becomes clearer.}

\emph{My beliefs soften and reshape themselves.}

The \textbf{Empathy Circle} has become one of my favorite places to explore this.\index{empathy circle}

It's a place where I can:

\begin{itemize}
\item speak honestly
\item hear myself think
\item and gently try on new beliefs
\end{itemize}

It's where I practice declaring:

\textbf{This is what I choose now.}

\emph{And little by little, I begin to act in a new way.}

\section{A Quiet Closing Thought}

I believe I can discover the beliefs that don't serve me . . . 

and change them.

Not all at once.\\
Not perfectly.

But steadily.

Kindly.

And that, to me, is one of the most beautiful better thoughts of all.
